% ==========================================================================
% network.tex  -  self-contained "Leader network load" subsection.
%
% Generated by SCRIPTS/aggregate_network.py from the leader-side traffic capture
% in SCRIPTS/network_data/ (iso_leader.py's TrafficMeter output). Re-run the
% script to refresh the tables.
%
% Compiles on its own (pdflatex network.tex). To fold it into the paper, drop the
% body between \begin{document} and \end{document} into the agent-scaling
% subsection next to the power measurement.
% ==========================================================================
\documentclass[11pt]{article}
\usepackage[a4paper,margin=2cm]{geometry}
\usepackage{booktabs}
\usepackage{array}
\usepackage{float}
\usepackage{caption}
\setlength{\tabcolsep}{5pt}

\begin{document}

\subsection{Leader network load}
\label{subsub:network:load}

This subsection quantifies the traffic the leader agent sustains as the fleet of
sensing agents grows from $1$ to $100$. A mock leader runs the live network
overlay (zero-configuration UDP discovery plus pull-on-join tool aggregation) with
no language model, election or dispatch, while a fleet of fake sensing agents is
brought up in waves of $1$, $10$, $20$, $50$ and $100$, each size held for one
minute. Every message the leader sends or receives is logged with its
application-level payload size. The reported bytes are therefore the JSON payload
on the wire only: they exclude the TCP/IP and Ethernet headers, the
acknowledgements and the per-message TCP handshake, which a connection-per-message
transport pays on top and which an on-wire capture would additionally count.

Table~\ref{tab:net_load} reports the per-minute message and byte rates the
leader sends and receives at each fleet size, plus a leader-idle baseline measured
after the run. Table~\ref{tab:net_kind} breaks the received bytes down by message
type. The received load is dominated by discovery: every agent re-announces its
presence every two seconds, so the HELLO beacons the leader takes in grow linearly
with the number of agents, at about 1.8~KiB/min per agent (least-squares
fit over the five sizes), reaching 184.8~KiB/min at $100$ agents. The one-time
{\ttfamily tools/list} aggregation (request out, result in) is a short burst at
each wave and is amortised to a small rate here; the leader's own announce and its
broadcast echo form the fixed idle floor.

\begin{table}[H]
\centering
\caption{Leader network load as the number of sensing agents grows. Rates are
averaged over the one-minute hold at each size (window length in the second
column); the idle row is the leader alone after the fleet stopped.
Application-level JSON payload only, excluding TCP/IP/Ethernet and handshake
overhead.}
\label{tab:net_load}
\begin{tabular}{cccccc}
\toprule
\textbf{Agents} & \textbf{Window [s]} & \textbf{Msgs/min in} & \textbf{Msgs/min out} & \textbf{In [KiB/min]} & \textbf{Out [KiB/min]} \\
\midrule
0 (idle) & 431 & 30 & 30 & 1.7 & 1.7 \\
1 & 60 & 76 & 40 & 5.9 & 2.6 \\
10 & 60 & 351 & 43 & 22.5 & 2.8 \\
20 & 60 & 642 & 62 & 36.8 & 4.5 \\
50 & 60 & 1578 & 81 & 96.2 & 6.2 \\
100 & 60 & 2833 & 130 & 184.8 & 10.5 \\
\bottomrule
\end{tabular}
\end{table}

\begin{table}[H]
\centering
\caption{Received bytes at the leader, per minute, split by message type: fleet
HELLO beacons (scale linearly with the agent count), {\ttfamily tools/list}
results (one-time aggregation burst per wave, amortised over the minute), and the
leader's own broadcast echo (fixed). KiB/min, application payload only.}
\label{tab:net_kind}
\begin{tabular}{ccccc}
\toprule
\textbf{Agents} & \textbf{HELLO (fleet)} & \textbf{tools/list result} & \textbf{HELLO (self-echo)} & \textbf{Total in} \\
\midrule
0 (idle) & 0.0 & 0.0 & 1.7 & 1.7 \\
1 & 2.2 & 2.0 & 1.7 & 5.9 \\
10 & 17.6 & 3.2 & 1.7 & 22.5 \\
20 & 34.6 & 0.6 & 1.7 & 36.8 \\
50 & 86.1 & 8.5 & 1.6 & 96.2 \\
100 & 152.9 & 30.2 & 1.7 & 184.8 \\
\bottomrule
\end{tabular}
\end{table}

\end{document}
